%%%%%%%%%%%%%%%%%%%%%%%%%%%%%%%%%%%%%%%%%%%%%%%%%%%%%%%%%%%%%%%%%%%%%%%%%%%%%%%
% THE LIMITS OF UTILITY MAXIMIZATION: EVIDENCE FOR COMPLEMENTARITY AND CONTEXT
%
% Based on: Chapter 3 of the Evidence-Based Framework (EBF)
% Style: Ernst Fehr (Zurich Tradition)
% Target: American Economic Review / Journal of Political Economy
%
% 8D Coordinates: D1=0.9, D2=0.9, D3=0.7, D4=0.8, D5=G1, D6=1.0, D7=0.1, D8=0.95
%%%%%%%%%%%%%%%%%%%%%%%%%%%%%%%%%%%%%%%%%%%%%%%%%%%%%%%%%%%%%%%%%%%%%%%%%%%%%%%

\documentclass[12pt]{article}

\usepackage{amsmath,amssymb,amsthm}
\usepackage{graphicx}
\usepackage{booktabs}
\usepackage{natbib}
\usepackage{setspace}
\usepackage[margin=1in]{geometry}
\usepackage{hyperref}

\newtheorem{proposition}{Proposition}
\newtheorem{hypothesis}{Hypothesis}
\newtheorem{prediction}{Prediction}
\newtheorem{definition}{Definition}

\doublespacing

\title{\textbf{The Limits of Utility Maximization under Stable Preferences:\\
Evidence for Complementarity and Context Dependence}}

\author{
Evidence-Based Framework Research Team\thanks{FehrAdvice \& Partners AG. Correspondence: research@fehradvice.com. We thank the behavioral economics community for decades of rigorous experimental evidence that makes this synthesis possible.}
}

\date{January 2026}

\begin{document}

\maketitle

%%%%%%%%%%%%%%%%%%%%%%%%%%%%%%%%%%%%%%%%%%%%%%%%%%%%%%%%%%%%%%%%%%%%%%%%%%%%%%%
% ABSTRACT
%%%%%%%%%%%%%%%%%%%%%%%%%%%%%%%%%%%%%%%%%%%%%%%%%%%%%%%%%%%%%%%%%%%%%%%%%%%%%%%

\begin{abstract}
\noindent
We examine the empirical and theoretical limits of utility maximization under stable preferences. Drawing on evidence from six domains---social preferences, identity and self-image, reference dependence, framing effects, preference construction, and temporal discounting---we demonstrate that the stability assumption is not merely a simplification but a substantive restriction that excludes central economic phenomena. The deeper problem is \textit{structural}: utilities exhibit interdependence across agents ($\partial^2 U_i / \partial x_i \partial x_j \neq 0$), time ($\partial U_t / \partial a_{t-1} \neq 0$), and context ($\partial U / \partial \Psi \neq 0$). These cross-derivatives are precisely the \textit{complementarities} that generate coordination problems, multiple equilibria, and path dependence. We formalize these interdependencies within a unified framework where complementarity $C$ and context $\Psi$ are explicit parameters. The framework does not reject utility maximization but embeds it within a richer structure that accommodates 50+ years of behavioral evidence. We derive testable predictions about when behavioral deviations are large versus small, and propose experimental designs to estimate context-dependent parameters.

\medskip
\noindent
\textbf{JEL Codes:} D01, D91, D87, C92\\
\textbf{Keywords:} Utility maximization, stable preferences, social preferences, reference dependence, complementarity, context dependence
\end{abstract}

\newpage

%%%%%%%%%%%%%%%%%%%%%%%%%%%%%%%%%%%%%%%%%%%%%%%%%%%%%%%%%%%%%%%%%%%%%%%%%%%%%%%
% 1. INTRODUCTION
%%%%%%%%%%%%%%%%%%%%%%%%%%%%%%%%%%%%%%%%%%%%%%%%%%%%%%%%%%%%%%%%%%%%%%%%%%%%%%%

\section{Introduction}
\label{sec:intro}

The classical framework in economics assumes that agents maximize well-defined utility functions with stable preferences. This assumption is not arbitrary---it serves crucial theoretical functions: tractability (fixed preferences enable optimization), prediction (stable preferences allow demand estimation), welfare analysis (unchanging preferences define a welfare criterion), and identification (without stability, any behavior can be rationalized as ``rational'').

Yet five decades of experimental and field evidence have accumulated demonstrating systematic violations of this assumption. The question is no longer \textit{whether} violations occur but \textit{what structure} underlies them. This paper argues that the violations share a common feature: \textbf{utility interdependence}. Specifically:
\begin{itemize}
    \item \textbf{Cross-agent interdependence:} $\frac{\partial U_i}{\partial x_j} \neq 0$ (social preferences)
    \item \textbf{Temporal interdependence:} $\frac{\partial U_t}{\partial a_{t-1}} \neq 0$ (habits, identity)
    \item \textbf{Intertemporal inconsistency:} $U_t(c_t, c_{t+1}) \neq \delta \cdot U_{t+1}(c_{t+1}, c_{t+2})$ (present bias)
    \item \textbf{Context interdependence:} $\frac{\partial U}{\partial \Psi} \neq 0$ (framing, defaults)
\end{itemize}

These interdependencies generate non-zero second derivatives---\textit{complementarities}---that are precisely what classical theory assumes away.

\subsection{The Stigler-Becker Defense and Its Limits}

\citet{stiglerbecker1977} argued that apparent preference changes can always be reinterpreted as responses to changing constraints: ``De gustibus non est disputandum''---tastes don't change; only opportunities do. For example, addiction can be modeled as a stable preference over consumption and a stock variable $S_t = \sum_{\tau < t} c_\tau$:
\begin{equation}
U(c_t, S_t) \quad \text{where } S_t \text{ is consumption stock}
\end{equation}

This reinterpretation saves stability but at a cost: it makes the theory unfalsifiable. Any behavior pattern can be rationalized with a sufficiently complex constraint structure. The question becomes empirical: do we gain predictive power from treating preferences as stable, or from explicitly modeling their dependence on social context, reference points, and time?

\subsection{Contribution}

We make three contributions. First, we synthesize evidence from six domains where the stable preferences assumption fails, showing that these are not isolated anomalies but manifestations of a common structure. Second, we formalize the underlying interdependencies within a framework where complementarity $C$ and context $\Psi$ are explicit parameters. Third, we derive testable predictions about when behavioral deviations are large versus small, enabling quantitative prediction rather than mere cataloguing of anomalies.


%%%%%%%%%%%%%%%%%%%%%%%%%%%%%%%%%%%%%%%%%%%%%%%%%%%%%%%%%%%%%%%%%%%%%%%%%%%%%%%
% 2. SIX DOMAINS WHERE STABLE PREFERENCES FAIL
%%%%%%%%%%%%%%%%%%%%%%%%%%%%%%%%%%%%%%%%%%%%%%%%%%%%%%%%%%%%%%%%%%%%%%%%%%%%%%%

\section{Evidence: Six Domains of Instability}
\label{sec:evidence}

\subsection{Domain 1: Social Preferences}

\textbf{The Evidence.} Thousands of experiments demonstrate that people care about others' outcomes:
\begin{itemize}
    \item \textbf{Ultimatum games:} Responders reject unfair offers, sacrificing money \citep{guth1982}
    \item \textbf{Dictator games:} Dictators share even when anonymous \citep{forsythe1994}
    \item \textbf{Public goods games:} Many contribute despite dominant strategy to free-ride \citep{andreoni1988}
    \item \textbf{Trust games:} Investors trust and trustees reciprocate \citep{berg1995}
\end{itemize}

The ultimatum game illustrates most starkly: a proposer divides \$10, and a responder accepts or rejects (rejection yields zero for both). Classical prediction: offer \$0.01, accept. Actual behavior: modal offers are \$4--5; offers below \$2 are frequently rejected.

\begin{proposition}[Social Preference Implication]
Utility is not $U_i(x_i)$ but $U_i(x_i, x_j)$. The cross-derivative is non-zero:
\begin{equation}
\frac{\partial^2 U_i}{\partial x_i \partial x_j} \neq 0
\end{equation}
This is precisely the complementarity that classical theory assumes away.
\end{proposition}

\textbf{Heterogeneity.} The distribution of types is well-established \citep{fehrcharness2024}:

\begin{table}[h]
\centering
\begin{tabular}{lcc}
\toprule
\textbf{Type} & \textbf{Students} & \textbf{General Population} \\
\midrule
Selfish & 40--60\% & 20--30\% \\
Inequality averse & 10--15\% & 15--20\% \\
Altruistic & 25--45\% & 45--55\% \\
\bottomrule
\end{tabular}
\caption{Distribution of social preference types}
\label{tab:types}
\end{table}

Critically, the \textit{majority} of individuals exhibit social preferences; purely selfish individuals are a minority, especially in representative samples.

\textbf{Context Dependence.} Social preferences vary with context:
\begin{itemize}
    \item Higher in face-to-face interactions
    \item Higher with in-group members
    \item Lower in competitive frames
    \item Sensitive to merit: inequality from effort is more accepted than inequality from luck \citep{cappelen2007,almas2020}
\end{itemize}

Thus the complementarity itself is context-dependent: $C_{ij} = C_{ij}(\Psi)$.

\subsection{Domain 2: Identity and Self-Image}

\textbf{The Evidence.} People sacrifice material payoffs to maintain consistent self-image. The ``moral wiggle room'' experiment \citep{danaberton2007} is instructive: subjects chose between (A) \$6 for self, \$1 for other, or (B) \$5 for self, \$5 for other. Many chose B (fair). But when given the option to ``not know'' what the other receives, many chose ignorance and then A. They preferred not knowing they were being unfair.

\begin{proposition}[Identity Utility]
Utility includes self-image:
\begin{equation}
U = U^{material} + U^{IDN}
\end{equation}
where $U^{IDN}$ is identity utility---the value of being (or seeing oneself as) a good person.
\end{proposition}

Identity creates intertemporal linkage:
\begin{equation}
\frac{\partial^2 U}{\partial a_t \partial a_{t+1}} > 0
\end{equation}
Behaving as type $X$ today makes it easier to see oneself as type $X$ tomorrow. This mechanism underlies habit, character, and culture.

\subsection{Domain 3: Reference Dependence and Loss Aversion}

\textbf{The Evidence.} \citet{kahnemantversky1979} demonstrated that outcomes are evaluated relative to a reference point, with losses weighted more than gains:
\begin{equation}
v(x) = \begin{cases}
x^\alpha & \text{if } x \geq 0 \\
-\lambda (-x)^\beta & \text{if } x < 0
\end{cases}
\end{equation}
with $\lambda \approx 2$ (loss aversion coefficient).

The endowment effect \citep{knetschetal1990} provides clean evidence: subjects given mugs or chocolate bars trade far below 50\%---they value what they have more than equivalent items they don't have.

\begin{proposition}[Reference Point Dependence]
The reference point $r$ is part of context:
\begin{equation}
U = U(x - r(\Psi))
\end{equation}
where $r(\Psi)$ depends on status quo, expectations \citep{koszegirab2006}, social comparison, and aspirations.
\end{proposition}

\subsection{Domain 4: Framing and Context Effects}

\textbf{The Evidence.} Identical choice problems elicit different responses depending on framing:
\begin{itemize}
    \item \textbf{Asian Disease Problem:} ``200 saved'' vs. ``400 die'' reverses preferences \citep{tverskykahneman1981}
    \item \textbf{Default effects:} Germany (opt-in): 12\% organ donation; Austria (opt-out): 99\% \citep{johnsonetal2003}
    \item \textbf{Anchoring:} Arbitrary numbers influence willingness to pay
\end{itemize}

\begin{proposition}[Frame Dependence]
Choices depend on presentation:
\begin{equation}
\text{Choice}(A, B; \Psi_1) \neq \text{Choice}(A, B; \Psi_2)
\end{equation}
even when $A$ and $B$ are objectively identical.
\end{proposition}

This directly violates stable preferences, which requires: same options $\Rightarrow$ same choice.

\subsection{Domain 5: Preference Construction}

\textbf{The Evidence.} For many choices, preferences are not ``retrieved'' but ``constructed'' at decision time. Contingent valuation studies show willingness to pay varies wildly with question order, starting points, and comparison goods.

\begin{proposition}[Constructed Preferences]
For complex goods, preferences emerge from deliberation in context:
\begin{equation}
U(x; \Psi) = f(\text{construction process}(\Psi))
\end{equation}
This is especially true for emotional, social, identity, and ecological utility.
\end{proposition}

\subsection{Domain 6: Temporal Preferences and Present Bias}

\textbf{The Evidence.} \citet{laibson1997} formalized quasi-hyperbolic discounting:
\begin{equation}
U_t = u(c_t) + \beta \sum_{\tau=1}^{T} \delta^\tau u(c_{t+\tau}) \quad \text{where } \beta < 1
\end{equation}

Empirical estimates consistently find $\beta \approx 0.7$ and $\delta \approx 0.99$ \citep{frederick2002,cohenetal2020}. This implies:
\begin{itemize}
    \item Substantial present bias (30\% premium on immediate rewards)
    \item Time-inconsistent behavior: plans made today are not followed tomorrow
    \item Demand for commitment devices \citep{ashraf2006,bryan2010}
\end{itemize}

\begin{proposition}[Temporal Context]
Temporal preferences constitute a $\Psi_T$ dimension of context:
\begin{equation}
\Psi_T = (\beta, \delta, \hat{\beta})
\end{equation}
where $\hat{\beta}$ reflects sophistication about own bias. This interacts with other context dimensions:
\begin{itemize}
    \item $\Psi_T \times \Psi_C$: Cognitive load worsens self-control
    \item $\Psi_T \times \Psi_I$: Institutions can provide commitment
\end{itemize}
\end{proposition}


%%%%%%%%%%%%%%%%%%%%%%%%%%%%%%%%%%%%%%%%%%%%%%%%%%%%%%%%%%%%%%%%%%%%%%%%%%%%%%%
% 3. THE DEEPER PROBLEM: UTILITY INTERDEPENDENCE
%%%%%%%%%%%%%%%%%%%%%%%%%%%%%%%%%%%%%%%%%%%%%%%%%%%%%%%%%%%%%%%%%%%%%%%%%%%%%%%

\section{The Deeper Problem: Structural Interdependence}
\label{sec:interdependence}

Each domain above could be treated as a ``behavioral anomaly''---a deviation requiring a patch. But the deeper issue is \textit{structural}: utilities are interdependent in ways that classical theory cannot accommodate.

\begin{definition}[Four Types of Interdependence]
\begin{enumerate}
    \item \textbf{Cross-agent:} $\frac{\partial U_i}{\partial x_j} \neq 0$ (my utility depends on your outcome)
    \item \textbf{Temporal:} $\frac{\partial U_t}{\partial a_{t-1}} \neq 0$ (current utility depends on past actions)
    \item \textbf{Intertemporal inconsistency:} $U_t(\cdot) \neq \delta \cdot U_{t+1}(\cdot)$ (time preferences not constant)
    \item \textbf{Context:} $\frac{\partial U}{\partial \Psi} \neq 0$ (utility depends on institutional environment)
\end{enumerate}
\end{definition}

\begin{proposition}[Complementarity Implication]
All four types generate non-zero second derivatives:
\begin{equation}
C_{ij} = \frac{\partial^2 U_i}{\partial x_i \partial x_j} \neq 0
\end{equation}
This is the structural feature that creates coordination problems, multiple equilibria, and path dependence.
\end{proposition}

\subsection{What Utility Maximization Gets Wrong}

The problem is not that utility maximization is false. Given a utility function and constraints, maximization is a reasonable model of choice. The problem is that the framework presupposes:
\begin{enumerate}
    \item The utility function is known and stable
    \item The constraints are exogenous
    \item Agents optimize independently
\end{enumerate}

When utilities are interdependent and context is endogenous, ``maximize utility'' doesn't define behavior---it merely redescribes whatever happens as ``utility-maximizing.''

\textbf{The Identification Problem.} If utility can include anything, any behavior is ``rational'' and the theory becomes unfalsifiable.

\textbf{The Equilibrium Selection Problem.} With interdependent utilities, multiple equilibria exist. Utility maximization says: at equilibrium, everyone is maximizing. It doesn't say which equilibrium will emerge.

\textbf{The Stability Problem.} Utility maximization describes optimal response to given conditions. It doesn't describe whether those conditions will persist.


%%%%%%%%%%%%%%%%%%%%%%%%%%%%%%%%%%%%%%%%%%%%%%%%%%%%%%%%%%%%%%%%%%%%%%%%%%%%%%%
% 4. THE REQUIRED EXTENSION
%%%%%%%%%%%%%%%%%%%%%%%%%%%%%%%%%%%%%%%%%%%%%%%%%%%%%%%%%%%%%%%%%%%%%%%%%%%%%%%

\section{A Framework with Complementarity and Context}
\label{sec:framework}

\subsection{From Optimization to Coherence}

The limits of utility maximization suggest a different criterion: not ``is each agent optimizing?'' but ``is the system of utilities coherent?''

\begin{definition}[Coherence]
A system is coherent if:
\begin{itemize}
    \item Expectations match outcomes
    \item Institutions support behavior
    \item Norms align with incentives
    \item Identity is consistent with action
\end{itemize}
\end{definition}

\subsection{The Extended Framework}

\begin{proposition}[Framework Extension]
The required extension includes:
\begin{enumerate}
    \item Allow $C_{ij} \neq 0$: Interdependent utilities with explicit complementarity parameter
    \item Allow $\Psi_{t+1} \neq \Psi_t$: Endogenous context that evolves
    \item Define $K$: Coherence as structural stability
    \item Define $Q$: Quality as welfare at equilibrium
    \item Analyze stability: When does $K \approx 1$ persist?
\end{enumerate}
\end{proposition}

The complementarity structure $C^*(\Psi)$ captures how utility interdependencies vary with context. This is not a minor extension---it transforms a single-agent optimization problem into a system problem.


%%%%%%%%%%%%%%%%%%%%%%%%%%%%%%%%%%%%%%%%%%%%%%%%%%%%%%%%%%%%%%%%%%%%%%%%%%%%%%%
% 5. TESTABLE PREDICTIONS
%%%%%%%%%%%%%%%%%%%%%%%%%%%%%%%%%%%%%%%%%%%%%%%%%%%%%%%%%%%%%%%%%%%%%%%%%%%%%%%

\section{Testable Predictions}
\label{sec:predictions}

The framework generates quantitative predictions about when behavioral deviations are large versus small.

\begin{prediction}[Social Preference Magnitude]
Social preference strength $\frac{\partial U_i}{\partial x_j}$ varies predictably:
\begin{equation}
\frac{\partial U_i}{\partial x_j} = f(\text{distance}_{ij}, \text{merit}, \text{frame}, \text{stakes})
\end{equation}
where distance is social/psychological distance, merit captures perceived deservingness, frame captures cooperative vs. competitive framing, and stakes captures decision magnitude.
\end{prediction}

\begin{hypothesis}[Context-Dependent Loss Aversion]
Loss aversion $\lambda$ varies with context:
\begin{equation}
\lambda(\Psi) = f(\Psi_C, \Psi_T, \Psi_E, \text{domain})
\end{equation}
where $\Psi_C$ is cognitive load, $\Psi_T$ is temporal context, and $\Psi_E$ is economic context.
\end{hypothesis}

\begin{prediction}[Present Bias Heterogeneity]
Present bias $\beta$ varies across:
\begin{itemize}
    \item Cognitive load: Higher load $\Rightarrow$ lower $\beta$ (more present-biased)
    \item Stakes: Small stakes show larger present bias
    \item Domain: Health behaviors show more present bias than financial decisions
\end{itemize}
\end{prediction}

\begin{prediction}[Framing Effect Magnitude]
Frame sensitivity varies inversely with preference crystallization:
\begin{equation}
|\text{Frame Effect}| \propto (1 - \text{Preference Familiarity})
\end{equation}
Novel decisions show larger framing effects; familiar decisions show smaller effects.
\end{prediction}


%%%%%%%%%%%%%%%%%%%%%%%%%%%%%%%%%%%%%%%%%%%%%%%%%%%%%%%%%%%%%%%%%%%%%%%%%%%%%%%
% 6. EXPERIMENTAL DESIGN
%%%%%%%%%%%%%%%%%%%%%%%%%%%%%%%%%%%%%%%%%%%%%%%%%%%%%%%%%%%%%%%%%%%%%%%%%%%%%%%

\section{Proposed Experimental Tests}
\label{sec:experiments}

\subsection{Experiment 1: Context-Dependent Social Preferences}

\textbf{Design:} $2 \times 2 \times 2$ factorial design varying:
\begin{enumerate}
    \item Frame (cooperative vs. competitive)
    \item Merit (earned vs. windfall endowment)
    \item Social distance (in-group vs. out-group)
\end{enumerate}

\textbf{Outcome:} Allocations in modified dictator game.

\textbf{Prediction:} Three-way interaction: social preferences highest in cooperative frame, earned endowment, in-group condition; lowest in competitive frame, windfall, out-group.

\subsection{Experiment 2: Loss Aversion Under Cognitive Load}

\textbf{Design:} Within-subject design measuring loss aversion under:
\begin{enumerate}
    \item No load (control)
    \item Working memory load (7-digit string)
    \item Time pressure
\end{enumerate}

\textbf{Outcome:} Accept/reject decisions for mixed gambles.

\textbf{Prediction:} Loss aversion increases under cognitive load: $\lambda_{load} > \lambda_{control}$.

\subsection{Experiment 3: Present Bias and Commitment}

\textbf{Design:} Longitudinal study with:
\begin{enumerate}
    \item Baseline $\beta$ elicitation
    \item Commitment device availability manipulation
    \item Follow-up behavior tracking
\end{enumerate}

\textbf{Prediction:} Individuals with lower $\beta$ (more present-biased) show higher demand for commitment devices and larger behavior improvement when devices are available.

\subsection{Identification Strategy}

Causal identification relies on:
\begin{enumerate}
    \item Random assignment to context conditions
    \item Within-subject designs where feasible
    \item Pre-registered analysis plans
    \item Multiple replications across populations
\end{enumerate}


%%%%%%%%%%%%%%%%%%%%%%%%%%%%%%%%%%%%%%%%%%%%%%%%%%%%%%%%%%%%%%%%%%%%%%%%%%%%%%%
% 7. FROM ANOMALIES TO ARCHITECTURE
%%%%%%%%%%%%%%%%%%%%%%%%%%%%%%%%%%%%%%%%%%%%%%%%%%%%%%%%%%%%%%%%%%%%%%%%%%%%%%%

\section{From Anomalies to Architecture}
\label{sec:architecture}

The behavioral economics literature spent decades documenting deviations from classical rationality. This was necessary and valuable---establishing \textit{that} people deviate. But the next step is to predict \textit{when} and \textit{how much}.

\begin{table}[h]
\centering
\small
\begin{tabular}{lll}
\toprule
\textbf{Domain} & \textbf{What Varies} & \textbf{Calibration Target} \\
\midrule
Social preferences & Cultures, stakes, merit & $C_{ij}(\Psi_S, \Psi_I)$ \\
Identity & Groups, priming, salience & $U^{IDN}(\Psi_S, \text{group})$ \\
Reference dependence & Domains, expectations & $r(\Psi)$, $\lambda(\Psi)$ \\
Framing & Presentations, defaults & Frame sensitivity $= f(\Psi_C)$ \\
Preference construction & Familiarity, complexity & Construction weight $= f(\Psi_K)$ \\
Temporal preferences & Cultures, cognitive load & $\beta(\Psi_C, \Psi_T)$ \\
\bottomrule
\end{tabular}
\caption{Calibration opportunities from behavioral domains}
\label{tab:calibration}
\end{table}

Each experimental finding that documents context-dependence adds information to the functional form of $C^*(\Psi)$. The variation that makes behavioral economics ``messy'' is precisely what makes calibration possible.


%%%%%%%%%%%%%%%%%%%%%%%%%%%%%%%%%%%%%%%%%%%%%%%%%%%%%%%%%%%%%%%%%%%%%%%%%%%%%%%
% 8. CONCLUSION
%%%%%%%%%%%%%%%%%%%%%%%%%%%%%%%%%%%%%%%%%%%%%%%%%%%%%%%%%%%%%%%%%%%%%%%%%%%%%%%

\section{Conclusion}
\label{sec:conclusion}

The limits of utility maximization under stable preferences are not anomalies to be patched but structural features to be modeled. Six domains---social preferences, identity, reference dependence, framing, preference construction, and temporal discounting---share a common feature: utility interdependence generating non-zero complementarities.

The required extension does not reject utility maximization but embeds it within a richer structure where complementarity $C$ and context $\Psi$ are explicit parameters. This transforms the question from ``what maximizes utility?'' to ``what pattern of interdependent utilities is coherent given context?''

The practical implications are substantial. Interventions that ignore context dependence will fail unpredictably. Policies designed for ``representative agents'' will miss the heterogeneity that determines aggregate outcomes. But the same features that make prediction difficult also make it possible: the variation across contexts is information, not noise.

Future work should calibrate context-dependent parameters from experimental variation and test the framework's predictions in field settings. The goal is not to abandon the optimization framework but to use it correctly---with the interdependencies and context effects that decades of evidence have documented.

%%%%%%%%%%%%%%%%%%%%%%%%%%%%%%%%%%%%%%%%%%%%%%%%%%%%%%%%%%%%%%%%%%%%%%%%%%%%%%%
% REFERENCES
%%%%%%%%%%%%%%%%%%%%%%%%%%%%%%%%%%%%%%%%%%%%%%%%%%%%%%%%%%%%%%%%%%%%%%%%%%%%%%%

\bibliographystyle{aer}

\begin{thebibliography}{99}

\bibitem[Almas et al.(2020)]{almas2020}
Almas, I., Cappelen, A. W., Sorensen, E. O., \& Tungodden, B. (2020).
Cutthroat capitalism versus cuddly socialism: Are Americans more meritocratic and efficiency-seeking than Scandinavians?
\textit{Journal of Political Economy}, 128(5), 1753--1788.

\bibitem[Andreoni(1988)]{andreoni1988}
Andreoni, J. (1988).
Why free ride? Strategies and learning in public goods experiments.
\textit{Journal of Public Economics}, 37(3), 291--304.

\bibitem[Ashraf et al.(2006)]{ashraf2006}
Ashraf, N., Karlan, D., \& Yin, W. (2006).
Tying Odysseus to the mast: Evidence from a commitment savings product in the Philippines.
\textit{Quarterly Journal of Economics}, 121(2), 635--672.

\bibitem[Berg et al.(1995)]{berg1995}
Berg, J., Dickhaut, J., \& McCabe, K. (1995).
Trust, reciprocity, and social history.
\textit{Games and Economic Behavior}, 10(1), 122--142.

\bibitem[Bryan et al.(2010)]{bryan2010}
Bryan, G., Karlan, D., \& Nelson, S. (2010).
Commitment devices.
\textit{Annual Review of Economics}, 2, 671--698.

\bibitem[Cappelen et al.(2007)]{cappelen2007}
Cappelen, A. W., Hole, A. D., Sorensen, E. O., \& Tungodden, B. (2007).
The pluralism of fairness ideals: An experimental approach.
\textit{American Economic Review}, 97(3), 818--827.

\bibitem[Cohen et al.(2020)]{cohenetal2020}
Cohen, J., Ericson, K. M., Laibson, D., \& White, J. M. (2020).
Measuring time preferences.
\textit{Journal of Economic Literature}, 58(2), 299--347.

\bibitem[Dana et al.(2007)]{danaberton2007}
Dana, J., Weber, R. A., \& Kuang, J. X. (2007).
Exploiting moral wiggle room: Experiments demonstrating an illusory preference for fairness.
\textit{Economic Theory}, 33(1), 67--80.

\bibitem[Fehr \& Charness(2024)]{fehrcharness2024}
Fehr, E., \& Charness, G. (2024).
Social preferences: Fundamental characteristics and economic consequences.
In \textit{Handbook of Experimental Economics} (Vol. 3). Princeton University Press.

\bibitem[Fehr \& Schmidt(1999)]{fehrschmidt1999}
Fehr, E., \& Schmidt, K. M. (1999).
A theory of fairness, competition, and cooperation.
\textit{Quarterly Journal of Economics}, 114(3), 817--868.

\bibitem[Forsythe et al.(1994)]{forsythe1994}
Forsythe, R., Horowitz, J. L., Savin, N. E., \& Sefton, M. (1994).
Fairness in simple bargaining experiments.
\textit{Games and Economic Behavior}, 6(3), 347--369.

\bibitem[Frederick et al.(2002)]{frederick2002}
Frederick, S., Loewenstein, G., \& O'Donoghue, T. (2002).
Time discounting and time preference: A critical review.
\textit{Journal of Economic Literature}, 40(2), 351--401.

\bibitem[G\"{u}th et al.(1982)]{guth1982}
G\"{u}th, W., Schmittberger, R., \& Schwarze, B. (1982).
An experimental analysis of ultimatum bargaining.
\textit{Journal of Economic Behavior \& Organization}, 3(4), 367--388.

\bibitem[Johnson et al.(2003)]{johnsonetal2003}
Johnson, E. J., \& Goldstein, D. (2003).
Do defaults save lives?
\textit{Science}, 302(5649), 1338--1339.

\bibitem[Kahneman \& Tversky(1979)]{kahnemantversky1979}
Kahneman, D., \& Tversky, A. (1979).
Prospect theory: An analysis of decision under risk.
\textit{Econometrica}, 47(2), 263--291.

\bibitem[Knetsch et al.(1990)]{knetschetal1990}
Kahneman, D., Knetsch, J. L., \& Thaler, R. H. (1990).
Experimental tests of the endowment effect and the Coase theorem.
\textit{Journal of Political Economy}, 98(6), 1325--1348.

\bibitem[K\H{o}szegi \& Rabin(2006)]{koszegirab2006}
K\H{o}szegi, B., \& Rabin, M. (2006).
A model of reference-dependent preferences.
\textit{Quarterly Journal of Economics}, 121(4), 1133--1165.

\bibitem[Laibson(1997)]{laibson1997}
Laibson, D. (1997).
Golden eggs and hyperbolic discounting.
\textit{Quarterly Journal of Economics}, 112(2), 443--478.

\bibitem[Stigler \& Becker(1977)]{stiglerbecker1977}
Stigler, G. J., \& Becker, G. S. (1977).
De gustibus non est disputandum.
\textit{American Economic Review}, 67(2), 76--90.

\bibitem[Tversky \& Kahneman(1981)]{tverskykahneman1981}
Tversky, A., \& Kahneman, D. (1981).
The framing of decisions and the psychology of choice.
\textit{Science}, 211(4481), 453--458.

\end{thebibliography}

%%%%%%%%%%%%%%%%%%%%%%%%%%%%%%%%%%%%%%%%%%%%%%%%%%%%%%%%%%%%%%%%%%%%%%%%%%%%%%%
% APPENDIX: 8D DOCUMENT TYPE SPECIFICATION
%%%%%%%%%%%%%%%%%%%%%%%%%%%%%%%%%%%%%%%%%%%%%%%%%%%%%%%%%%%%%%%%%%%%%%%%%%%%%%%

\appendix
\section{8D Document Type Specification}
\label{app:8d}

This paper was generated using the 8D Document Type Framework (Appendix DDD/CCC of the Evidence-Based Framework).

\textbf{Style Profile:} Ernst Fehr Academic Style

\textbf{8D Coordinates:}
\begin{itemize}
    \item $D_1$ (Knowledge): 0.9 --- Expert audience (behavioral economists, decision theorists)
    \item $D_2$ (Proximity): 0.9 --- Same field (behavioral economics)
    \item $D_3$ (Scope): 0.7 --- Field-wide impact
    \item $D_4$ (Time): 0.8 --- Full journal article reading time
    \item $D_5$ (Goal): G1 --- Inform (synthesize evidence, develop framework)
    \item $D_6$ (Context): 1.0 --- External/public (journal publication)
    \item $D_7$ (Emotion): 0.1 --- Highly factual
    \item $D_8$ (Persistence): 0.95 --- Archival
\end{itemize}

\textbf{Derived Structure (DT-5/DT-6):}
\begin{itemize}
    \item Technical glossary: Not required (expert audience)
    \item Detailed subsections: Yes (D4 $>$ 0.6)
    \item Theory section: Yes (D5 = G1)
    \item Full references: Yes (D8 $>$ 0.8)
\end{itemize}

\textbf{Derived Style (DT-7/DT-8):}
\begin{itemize}
    \item Formal register: Yes (D6 $>$ 0.7)
    \item Hedging required: Yes (D8 $>$ 0.8)
    \item Technical vocabulary: Yes (D1 $>$ 0.7)
    \item Impersonal style: Yes (D7 $<$ 0.3)
\end{itemize}

\section{Relationship to EBF}
\label{app:ebf}

This working paper synthesizes evidence presented in Chapter 3 of the Evidence-Based Framework for Economic and Social Behavior (EBF). The EBF is an integrative framework organizing 50+ years of behavioral economics evidence under the unified concepts of complementarity ($C$) and context ($\Psi$).

For complete formal derivations showing how each behavioral phenomenon emerges as a limiting case, see Appendix A of the EBF documentation.

\end{document}
